\section{META im BEATRIX-System}
%==============================================================================

\subsection{Rolle bei der LLM-basierten Estimation}

Das \BEATRIX-System verwendet ein LLM zur Parameterschätzung. Das \META-Modul informiert diesen Prozess auf mehreren Ebenen:

\begin{enumerate}
    \item \textbf{Base Rate Awareness (MA-EPI-26)}: Das LLM muss empirische Base Rates für Typen-Verteilungen berücksichtigen
    \item \textbf{Social Desirability Korrektur (MA-EPI-32)}: Persona-Beschreibungen müssen auf systematische Verzerrungen geprüft werden
    \item \textbf{Anchoring-Prävention (MA-EPI-14)}: Few-Shot Anchoring im Prompt verhindert willkürliche Schätzungen
\end{enumerate}

\subsection{Constrained Generation Pipeline}

Die \BEATRIX{} Estimation Pipeline (ESTIMATION-Modul) implementiert die epistemologischen Constraints:

\begin{enumerate}
    \item \textbf{PreFlight}: Prüfung auf plausibles Szenario (verhindert Halluzination)
    \item \textbf{Few-Shot Anchoring}: Kalibrierungsskalen im Prompt
    \item \textbf{Pydantic Schema}: Strukturierte Outputs
    \item \textbf{Ensemble Voting}: N parallele Runs, Median-Aggregation
    \item \textbf{IRR-Check}: Inter-Rater Reliability > 0.85
\end{enumerate}

\subsection{Ambiguity Detection}

Die Unterscheidung zwischen \textbf{Unsicherheit} (unimodale Verteilung) und \textbf{Wertekonflikt} (bimodale Verteilung) basiert auf:

\begin{equation}
BC = \frac{s^2 + 1}{k + 3 \cdot \frac{(n-1)^2}{(n-2)(n-3)}}
\end{equation}

wobei $s$ = Schiefe, $k$ = Kurtosis. Bei $BC > 0.555$ liegt Polarisierung vor.

%==============================================================================
